\documentclass[12pt]{article}

\usepackage[utf8]{inputenc}
\usepackage[T1]{fontenc}
\usepackage{geometry}
\geometry{a4paper, margin=1in}
\usepackage{graphicx}
\usepackage{hyperref}

\title{Manual for Creating a Document}
\author{Author Name}
\date{\today}

\begin{document}

\maketitle

\tableofcontents
\newpage

\section{Introduction}
This manual provides instructions on how to create a basic document using LaTeX.

\section{Document Structure}
A basic LaTeX document consists of a preamble and a document body.

\subsection{Preamble}
The preamble contains document class declarations, packages, and metadata like title, author, and date.

\subsection{Document Body}
The document body is enclosed within \texttt{\textbackslash begin\{document\}} and \texttt{\textbackslash end\{document\}} tags.

\section{Basic Formatting}
\begin{itemize}
    \item \textbf{Bold text} is created using \texttt{\textbackslash textbf\{\}}.
    \item \textit{Italic text} is created using \texttt{\textbackslash textit\{\}}.
    \item \underline{Underlined text} is created using \texttt{\textbackslash underline\{\}}.
\end{itemize}

\section{Lists}
You can create bulleted or numbered lists.

\subsection{Bulleted Lists}
\begin{itemize}
    \item Item 1
    \item Item 2
\end{itemize}

\subsection{Numbered Lists}
\begin{enumerate}
    \item First item
    \item Second item
\end{enumerate}

\section{Conclusion}
You now have the basics to start creating your own documents in LaTeX.

\end{document}
